\section{Maraschino: Cherry-Counts for TKF92}
\label{sec:maraschino-main}

We earlier introduced Maraschino as a TKF92 generalization of CherryML
in Section~\ref{sec:maraschino-intro}.  We now develop it in full.

The CherryML composite likelihood replaces the joint phylogenetic
likelihood with a product of pairwise sibling (``cherry'') likelihoods
\cite{Prillo2022CherryML}, exploiting the fact that for a fixed alignment,
the alignment-substitution counts are sufficient statistics for the
substitution rate parameters.  This reduces the training data---which can
be many gigabytes of alignments---to a fixed-shape counts tensor whose
size depends only on the alphabet $\alphabet$ and the number of
discretized cherry-time bins, and is independent of the number of
training alignments.  The resulting compact tensor fits in GPU memory
and may be optimized by standard autograd-based gradient methods.

\textbf{Maraschino} extends CherryML to TKF92 by additionally including
the alignment \emph{transition} counts as sufficient statistics for the
indel rates and fragment-extension probability.  Concretely, for each
pair of consecutive (non-empty) alignment columns we record the
source column type ($\xstate \in \{\sta,\mat,\ins,\del\}$),
the destination column type ($\ystate \in \{\mat,\ins,\del,\fin\}$),
and the ancestor/descendant characters in each column. We refer to
the resulting tensor of counts---binned by discretized pairwise
divergence time and aggregated over training cherries---as the
\emph{cherry-count tensor}, and the procedure for training TKF92 from
it as Maraschino.

The Maraschino composite log-likelihood for TKF92 is the sum, over
divergence-time bins~$b$ and adjacency types $\xstate \to \ystate$, of
the cherry-count multiplied by the model log-probability of that
adjacency at that bin's representative time:
\begin{multline}
\mathcal{L}_{\text{cherry}}^{\mathrm{TKF92}}(\insrate,\delrate,\ext,\exch,\eqm)
=
\sum_{b=1}^{n_\tau}\Bigg[
\sum_{\xstate,\ystate}
n^{(b)}_{\xstate\to\ystate}\,
\log \tkftrans'_{\xstate\to\ystate}(\insrate,\delrate,\bar\evoltime_b,\ext)
\\
+
\sum_{\anctok,\destok} n^{\mat,(b)}_{\anctok,\destok}
\log\big[ \eqm_\anctok \exp(\revsub \bar\evoltime_b)_{\anctok\destok} \big]
+
\sum_\destok n^{\ins,(b)}_\destok \log \eqm_\destok
+
\sum_\anctok n^{\del,(b)}_\anctok \log \eqm_\anctok
\Bigg]
\label{eq:maraschino-tkf92-ll}
\end{multline}
where $n^{(b)}_{\xstate\to\ystate}$ are the binned transition counts,
$n^{\mat,(b)}_{\anctok,\destok}, n^{\ins,(b)}_\destok, n^{\del,(b)}_\anctok$
are the binned per-state emission counts,
$\tkftrans'$ is the TKF92 Pair HMM transition matrix from
Section~\ref{sec:tkf92-machines},
and $\bar\evoltime_b$ is the representative time of bin~$b$
(typically the geometric mean of the bin endpoints).
The first term scores indel/fragment behaviour;
the remaining three score the substitution model.
Boundary contributions ($\sta\to\cdot$, $\cdot\to\fin$, $\sta\to\fin$
for empty alignments) are scored against the corresponding
$\sta$-row and $\fin$-column entries of $\tkftrans'$.

\paragraph{Why this is a sufficient-statistic compression.}
For fixed alignment and fixed discretized cherry times, the entire
collection of training cherries enters
equation~(\ref{eq:maraschino-tkf92-ll}) only through the binned
counts.  In particular, two cherry collections with identical counts
yield identical $\mathcal{L}_{\text{cherry}}^{\mathrm{TKF92}}$,
identical gradients, and identical maximum-likelihood estimates.
The counts tensor is therefore a sufficient statistic for the
TKF92 parameters under the cherry-count composite likelihood.

\paragraph{Sharded count accumulation.}
For a corpus such as Pfam-Full, cherries are extracted from each MSA
in parallel and contribute additively to the global counts tensor.
A simple map-reduce over Pfam clans accumulates the counts in
constant memory per worker; the global tensor is the
elementwise sum of per-shard tensors and is independent of the order
of accumulation.

\subsection{Optimizing Maraschino: gradient methods vs.\ inner EM}
\label{sec:maraschino-opt}

The cherry-count log-likelihood
$\mathcal{L}_{\text{cherry}}^{\mathrm{TKF92}}$ is a smooth function of
the unconstrained parameters
$(\log\insrate, \log\delrate, \log\ext, \log\exch, \log\eqm)$
(with the simplex and exchangeability constraints handled by softmax /
log-Cholesky reparameterisation).  Two distinct optimizers apply.

\paragraph{Maraschino-as-autograd.}
The objective is differentiable, so Adam, L-BFGS, and similar
optimizers may be applied directly via autograd.  Each step requires one matrix
exponential per cherry-time bin, and a forward+backward pass through
equation~(\ref{eq:maraschino-tkf92-ll}); both steps are highly
parallel on GPU.

\paragraph{Maraschino-around-EM.}
\label{sec:maraschino-around-em}
Alternatively, the full-batch EM iterates of Section~\ref{sec:bw-tkf92}
apply, with the cherry-count tensor playing the role of the E-step
output.  Concretely:
\begin{enumerate}[nosep]
\item For each bin~$b$, run a one-step expected-count update on the
  TKF92 Pair HMM at parameters $\theta^{(t)}$ and time $\bar\evoltime_b$,
  using the cherry-count tensor as the input ``observed
  alignment.''  This replaces the usual Forward--Backward E-step with a
  closed-form linear pass that resolves the $\mat,\ins,\del$ self-loop
  counts into fragment-extension counts and link-level counts
  (Section~\ref{sec:bw-tkf92}).
\item Accumulate the BDI sufficient statistics
  $(B,D,L,M,S,T)$ via the score identity
  (Section~\ref{sec:bridge-expectations}), and the bridge-expectation substitution
  statistics $(W,U,V)$ via the matched-pair contribution.
\item Apply the closed-form M-step
  (equations~(\ref{eq:kappa-quadratic})--(\ref{eq:insrate-root})
  for $\insrate,\delrate$;
  $\ext \leftarrow F/(F+E)$ for fragment extension;
  Section~\ref{sec:bw-tkf91} item~1 for $\exch,\eqm$).
\end{enumerate}
Because every M-step in Section~\ref{sec:bw-tkf92} is closed-form and
linear in the sufficient statistics, the inner EM iterates are exact
ascent on $\mathcal{L}_{\text{cherry}}^{\mathrm{TKF92}}$, with no
gradient-step or learning-rate hyperparameter.  In practice, EM
converges in ${\sim}10$--$30$ iterations on protein cherry-count
tensors and is competitive with Adam on the same objective.

\paragraph{Practical comparison.}
Adam handles arbitrary regularizers and is straightforward to
combine with stochastic minibatching over time-bins or alphabet
slices.  Inner EM does not require a learning rate, has monotonic
convergence guarantees, and produces the natural-gradient direction
``for free.''  When the cherry-count tensor fits in memory (it does,
even for Pfam-Full at $|\alphabet|=20$ and $n_\tau=20$), inner-EM
Maraschino is the recommended default.

\subsection{EM-around-Maraschino: mixtures of TKF92}
\label{sec:em-around-maraschino}

A mixture of $K$ TKF92 components, each with its own indel rates,
fragment-extension probability, and substitution model, captures
heterogeneity within and across protein families.  We fit such a
mixture by an outer EM loop with Maraschino as the inner optimizer
of each component:
\begin{algorithm}
\caption{EM-around-Maraschino for a mixture of $K$ TKF92 components.}
\label{alg:em-around-maraschino}
\begin{algorithmic}[1]
\State Initialize component parameters $\theta^{(0)}_1,\ldots,\theta^{(0)}_K$
       and mixture weights $\pi^{(0)}_1,\ldots,\pi^{(0)}_K$.
\For{$t = 0,1,\ldots$}
  \State \textbf{E-step.} For each cherry $c$ in the training set,
         compute the responsibility $\gamma^{(t)}_{c,k}
         \propto \pi^{(t)}_k\,P(c\mid \theta^{(t)}_k)$
         (the per-component cherry-count likelihood normalised across
         components).  Each $P(c\mid\theta_k)$ is a single forward pass
         through the cherry-count likelihood at the cherry's
         counts subtensor.
  \State \textbf{Component-weighted count accumulation.}
         For each component~$k$, accumulate the
         responsibility-weighted cherry-count tensor
         $C_k^{(t+1)} = \sum_c \gamma^{(t)}_{c,k}\,C_c$.
  \State \textbf{Component M-step (Maraschino).}
         For each component~$k$, run Maraschino
         (Section~\ref{sec:maraschino-around-em}) on $C_k^{(t+1)}$
         to obtain $\theta^{(t+1)}_k$.
  \State \textbf{Mixture-weight M-step.}
         $\pi^{(t+1)}_k \propto \sum_c \gamma^{(t)}_{c,k}$.
\EndFor
\end{algorithmic}
\end{algorithm}

\noindent
The outer EM is monotone in the mixture log-likelihood
$\sum_c \log\sum_k \pi_k P(c\mid\theta_k)$ because both the
responsibilities (E-step) and the per-component Maraschino solves
(M-step) are exact ascent on their respective Q-functions.

The decisive feature is that the outer M-step does \emph{not} require
re-extracting cherries or recomputing alignments per iteration:
the per-component cherry-count tensor $C_k^{(t+1)}$ is a
responsibility-weighted average of the same per-cherry tensors that
were extracted once and for all at the start of training.

\subsection{EM-around-CherryML: mixtures over site classes}
\label{sec:em-around-cherryml}

A complementary class of mixture model places the latent variable on
\emph{MSA columns} rather than on whole cherries: each MSA column~$c$
draws an unobserved \emph{site class}
$z_c \sim \catdist(\eqm_1,\ldots,\eqm_C)$ from a shared per-class
substitution model $\subproc(\exch^{(k)}, \eqm^{(k)})$, and every
cherry in the family that has a residue pair at column $c$ evolves
its match emission under the same class.  Because no indel rates
vary across mixture components, this is strictly a substitution-only
mixture, and the corresponding fitter is CherryML rather than
Maraschino: the cherry-count tensor of
Section~\ref{sec:maraschino-main} contributes only the matched
substitution observations $(\anctok, \destok, t_p)$ and the
indel-free composite likelihood reduces to that of
\cite{Prillo2022CherryML}.

\paragraph{Column-shared responsibilities.}
The defining observation is that within a single MSA, all cherries
that have a match at column $c$ inherit the \emph{same} latent class
$z_c$.  This couples cherries across the same column: the column-level
posterior
\[
\gamma_c(k)
\;\propto\;
\eqm_k \prod_{p \,:\, \mat \text{ at } c}
P\!\bigl(\anctok_{p,c}, \destok_{p,c} \mid k, t_p\bigr),
\qquad
P(\anctok, \destok \mid k, t)
= \eqm^{(k)}_\anctok \exp(\revsub^{(k)} t)_{\anctok\destok},
\]
multiplies the per-cherry match-emission likelihoods over all cherries
$p$ that contribute a residue pair to column $c$.  Cherries that have
a gap at $c$ contribute no factor for that column.  The resulting
composite log-likelihood
\begin{equation}\label{eq:em-cherryml-ll}
\mathcal{L}_{\text{site-mix}}
\;=\;
\sum_m \sum_{c \,\in\, \text{cols}(m)}
\log \sum_{k=1}^{C} \eqm_k
\prod_{p \,:\, \mat \text{ at } c}
\eqm^{(k)}_{\anctok_{p,c}}
\exp(\revsub^{(k)} t_p)_{\anctok_{p,c}\,\destok_{p,c}}
\end{equation}
is a sum, over MSAs $m$ and over columns $c$ within each MSA, of a
log-mixture in which cherries share the column's class.  This is in
direct contrast to EM-around-Maraschino
(Section~\ref{sec:em-around-maraschino}), whose responsibilities
$\gamma^{(t)}_{c,k}$ are per \emph{cherry}: there each cherry is
assigned to a single component independently of the columns it
spans.

\paragraph{Outer / inner EM.}
We fit~\eqref{eq:em-cherryml-ll} by alternating an MSA-wide
column-responsibility update with a per-class bridge-expectation M-step.
At iteration~$t$, the column responsibilities are
\begin{equation}\label{eq:em-cherryml-resp}
\gamma^{(t)}_{m,c}(k)
\;\propto\;
\eqm^{(t)}_k
\prod_{p\,:\, \mat \text{ at } c}
\eqm^{(t,k)}_{\anctok_{p,c}}\,
\exp\!\bigl(\revsub^{(t,k)}\, t_p\bigr)_{\anctok_{p,c}\,\destok_{p,c}},
\end{equation}
normalised across $k$ for each $(m, c)$, and the
responsibility-weighted bridge-expectation sufficient statistics for class
$k$ are
\begin{equation}\label{eq:em-cherryml-WUk}
\begin{aligned}
\hat W_i^{(t+1,k)}
  &= \sum_{m,c} \gamma^{(t)}_{m,c}(k)
     \sum_{p\,:\, \mat \text{ at } c}
     \emcount^W_i\!\bigl(\anctok_{p,c}, \destok_{p,c}, t_p\bigr),
\\
\hat U_{ij}^{(t+1,k)}
  &= \sum_{m,c} \gamma^{(t)}_{m,c}(k)
     \sum_{p\,:\, \mat \text{ at } c}
     \emcount^U_{ij}\!\bigl(\anctok_{p,c}, \destok_{p,c}, t_p\bigr),
\\
\hat V_i^{(t+1,k)}
  &= \sum_{m,c} \gamma^{(t)}_{m,c}(k)
     \sum_{p\,:\, \mat \text{ at } c}
     \delta\bigl(\anctok_{p,c} = i\bigr),
\end{aligned}
\end{equation}
where $\emcount^W_i, \emcount^U_{ij}$ are the standard
endpoint-conditioned bridge expectations
(Section~\ref{sec:bridge-expectations},
eqs.~\eqref{eq:em-dwell}--\eqref{eq:em-counts}) at the class-$k$
rate matrix $\revsub^{(t,k)}$.  The full outer iteration is then
Algorithm~\ref{alg:em-around-cherryml}.
\begin{algorithm}
\caption{EM-around-CherryML for $C$ site classes.}
\label{alg:em-around-cherryml}
\begin{algorithmic}[1]
\State Initialize per-class GTR parameters $\theta^{(0)}_k$ and class
weights $\eqm^{(0)}_k$.
\For{$t = 0, 1, \ldots$}
  \State \textbf{Outer E-step.} Compute column responsibilities
  $\gamma^{(t)}_{m,c}(k)$ via~\eqref{eq:em-cherryml-resp}.
  \State \textbf{Sufficient-statistic accumulation.} Compute
  $(\hat W^{(t+1,k)}, \hat U^{(t+1,k)}, \hat V^{(t+1,k)})$
  via~\eqref{eq:em-cherryml-WUk}.
  \State \textbf{Per-class GTR inner EM.} For each class $k$, run the
  closed-form GTR coordinate-ascent M-step of
  Section~\ref{sec:bw-tkf91} on
  $(\hat W^{(t+1,k)}, \hat U^{(t+1,k)}, \hat V^{(t+1,k)})$ to obtain
  $\theta^{(t+1)}_k$.  Because the responsibility weights are fixed,
  this step is a standard CherryML rate fit on a single GTR model
  with weighted observations.
  \State \textbf{Class-weight update.}
  $\eqm^{(t+1)}_k \propto \sum_{m,c} \gamma^{(t)}_{m,c}(k)$.
\EndFor
\end{algorithmic}
\end{algorithm}

\noindent
The outer iteration is monotone in
$\mathcal{L}_{\text{site-mix}}$: the column responsibilities are the
exact posterior over $z_c$ and the per-class GTR M-step is exact
ascent on its (responsibility-weighted) Q-function.  The inner GTR
EM is itself iterative because the GTR M-step couples
$(\exch^{(k)}, \eqm^{(k)})$ nonlinearly
(Section~\ref{sec:bw-tkf91}, item~2 of the M-step list); a few
inner sweeps per outer iteration suffice in practice.

\paragraph{Why CherryML, not Maraschino.}
No indel rates appear in~\eqref{eq:em-cherryml-ll}: the model assumes
the alignment is given and that all alignment columns share a single
TKF92 indel process whose parameters are held fixed (or estimated
elsewhere, e.g.\ by Maraschino on the same cherries).  Mixing
\emph{only} the substitution model means the
$(\insrate, \delrate, \ext)$-related row-normalisation and
fragment-extension responsibility issues of TKF92 do not arise;
the relevant sufficient statistics are exactly those of CherryML.
The only structural change relative to plain CherryML is the per-MSA,
per-column responsibility weighting in~\eqref{eq:em-cherryml-WUk},
which couples cherries that share columns of the same MSA.

\paragraph{Relation to per-class GTR mixtures elsewhere.}
The per-class GTR M-step of equation~\eqref{eq:em-cherryml-WUk} is
exactly the reversible-mixture update of
Appendix~\ref{sec:mstep-mixture}; the only difference between this
appendix and Algorithm~\ref{alg:em-around-cherryml} is the
provenance of the responsibilities (here: cross-cherry column
sharing within an MSA; in the appendix: any user-supplied per-site
posterior).  Site-mixture parameters fitted by
EM-around-CherryML can therefore be plugged into
mixture-aware downstream methods --- including Maraschino with a
fixed indel-process backbone --- without further re-derivation.

\subsection{Expected sufficient statistics as fast custom VJPs}
\label{sec:suffstat-vjp}

The Maraschino objective
(equation~(\ref{eq:maraschino-tkf92-ll})) is concretely a sum of
log-probabilities, each of which has a closed-form gradient via the
score identity and the BDI / bridge-expectation sufficient statistics.
Specifically, the $(\insrate,\delrate)$-gradient of
$\log \tkftrans'_{\xstate\to\ystate}(\insrate,\delrate,t,\ext)$
is given exactly by the score-derivative tables
of Sections~\ref{sec:score-table-general}--\ref{sec:score-table-limits},
and the $(\exch,\eqm)$-gradient of
$\log[\eqm_\anctok \exp(\revsub t)_{\anctok\destok}]$ is given by the
bridge-expectation dwell times and jump counts
(equations~(\ref{eq:em-dwell})--(\ref{eq:em-counts})).

For training pipelines built around stochastic-gradient optimizers
such as Adam, this observation enables a substantial speedup.
Rather than relying on autograd to back-propagate through a matrix
exponential and an eigendecomposition, we register a
\emph{custom vector-Jacobian product (VJP)} that returns the
score-derivative-and-bridge-expectation closed forms directly:
\begin{itemize}[nosep]
\item Forward pass: compute and cache the matrix exponential
$P(t) = e^{\revsub t}$ and the BDI link-level transition factors
$(\alpha,\beta,\gamma)$ at each cherry-time bin.
\item Backward pass: contract the upstream cotangent with the
expected sufficient statistics (bridge expectations and BDI score
expressions) instead of with the autograd-computed Jacobian.
\end{itemize}
The forward and backward costs are both
$O(|\alphabet|^2)$ per bin per emission state for the substitution
part and $O(1)$ for the indel part, matching the cost of evaluating
the closed-form M-step itself.

This pattern generalises beyond Maraschino.  Any model whose
M-step is closed-form linear in expected sufficient statistics admits
a corresponding fast custom VJP that exposes those expectations as the
gradient of the observed log-likelihood (this is, again, the score
identity~(\ref{eq:score-identity})).  Concretely, this enables Adam
and other gradient-based optimizers to train TKF91, TKF92, and
mixture-of-TKF92 models at a per-step cost matching exact EM, while
retaining Adam's flexibility to incorporate arbitrary differentiable
regularizers (e.g., weight decay on the BDI rates, KL-divergence
priors on the GTR exchangeability), to interleave updates with
non-EM-amenable parameters (e.g., the per-component mixture
weights of Section~\ref{sec:em-around-maraschino} when fit
jointly with neural-network embeddings), and to work
with the mini-batch SVI-BW pseudocount stream of
Section~\ref{sec:svi-bw}.
